\subsection{CONTINUOUS HOMOMORPHISMS AND STRICT MORPHISMS}

PROPOSITION 23. A homomorphism \(f\) of a topological group \(\mathbf{G}\) into a topological group \(\mathbf{G}'\) is continuous on \(\mathbf{G}\) if and only if it is continuous at one point of \(\mathbf{G}\) .

Suppose \(f\) is continuous at a point \(a \in \mathbf{G}\) ; then if \(\mathbf{V}'\) is any neighbourhood of \(f(a)\) , \(\mathbf{V} = f^{-1}(\mathbf{V}')\) is a neighbourhood of \(a\) . Hence if \(x\) is any point of \(\mathbf{G}\) , we have

\[
f (x a ^ {- 1} \mathrm{V}) = f (x) [ f (a) ] ^ {- 1} f (\mathrm{V}) \subset f (x) [ f (a) ] ^ {- 1} \mathrm{V} ^ {\prime},
\]

and therefore \(f\) is continuous at \(x\) .

A continuous homomorphism of a topological group G into a topological group \(G'\) is also called a morphism of G into \(G'\) for the topological group structures.

Let \(f\) be a continuous homomorphism of a topological group \(G\) into a topological group \(G'\) ; the inverse image \(H = f^{-1}(e')\) of the identity element \(e'\) of \(G'\) is a normal subgroup of \(G\) , and \(f(G)\) is a subgroup of \(G'\) . Consider the canonical factorization \(f = \psi \circ \overline{f} \circ \varphi\) , where \(\varphi\) is the canonical mapping \(G \to G/H\) , \(\psi\) is the canonical injection \(f(G) \to G'\) and lastly \(\overline{f}\) is a continuous bijective homomorphism of the quotient group \(G/H\) onto the subgroup \(f(G)\) (Chapter I, § 3, no. 5); \(\overline{f}\) is said to be the bijective homomorphism associated with \(f\) . In general, \(f\) is not an isomorphism of topological groups.

For example, let \(G'\) be a non-discrete topological group, and G the topological group obtained by giving \(G'\) the discrete topology; then the identity mapping of G into \(G'\) is a continuous bijective homomorphism, but is not bicontinuous.

DEFINITION 1. A continuous homomorphism of a topological group G into a topological group \(\mathbf{G}'\) is said to be a strict morphism of G into \(\mathbf{G}'\) if the bijective homomorphism \(\overline{f}\) of \(\mathbf{G}/f^{-1}(e')\) onto \(f(\mathbf{G})\) , associated with \(f\) , is an isomorphism of topological groups (in other words, if \(\overline{f}\) is bicontinuous).

An isomorphism of a topological group G onto a topological group \(G'\) is therefore a bijective strict morphism of G onto \(G'\) .

PROPOSITION 24. Let \(f\) be a continuous homomorphism of a topological group \(G\) into a topological group \(G'\) . Then the following three statements are equivalent: a) \(f\) is a strict morphism.

\begin{enumerate}
\def\labelenumi{\alph{enumi})}
\setcounter{enumi}{1}
\item
  The image under \(f\) of every open set in \(\mathbf{G}\) is an open set in \(f(\mathbf{G})\) .
\item
  The image under \(f\) of every neighbourhood of the identity element in \(G\) is a neighbourhood of the identity element in \(G'\) .
\end{enumerate}

In view of Lemma 2 of no. 4, the equivalence of a) and b) follows immediately from the definitions (Chapter I, § 5, no. 3, Proposition 5). The equivalence of b) and c) is a particular case of Proposition 15 of no. 5, if we observe that G operates continuously on \(f(\mathbf{G})\) by the external law \((s, f(t)) \to f(st)\) .

Remarks. 1) From condition b) of Proposition 24 it follows that every continuous homomorphism of a topological group into a discrete group is a strict morphism.

If G is compact and \(f(\mathbf{G})\) is Hausdorff, then the bijective homomorphism \(\overline{f}\) associated with f is bicontinuous (Chapter I, § 10, no. 2, Theorem 1, Corollary 2, and no. 1, Proposition 5, Corollary 4). Hence every continuous homomorphism of a compact group into a Hausdorff group is a strict morphism.

\pandocbounded{\includegraphics[keepaspectratio]{images/part002_802da4a7.jpg}}

\begin{enumerate}
\def\labelenumi{\arabic{enumi})}
\setcounter{enumi}{1}
\item
  Let \(f\) be a strict morphism of \(G\) into \(G'\) and let \(g\) be a strict morphism of \(G'\) into \(G''\) . If \(f\) is surjective or if \(g\) is injective, it follows immediately from Proposition 24 that \(g \circ f\) is a strict morphism of \(G\) into \(G''\) . But this conclusion is no longer necessarily valid if neither of the two preceding conditions is satisfied, even if \(f\) is injective and \(g\) surjective (Exercise 19).
\item
  Let \(f\) be a continuous homomorphism of a topological group \(G\) into a topological group \(G'\) , and let \(H\) be a normal subgroup of \(G\) . \(f\) induces a homomorphism \(g\) of the group \(G / H\) onto the quotient group \(f(G) / f(H)\) . This homomorphism \(g\) is continuous. Moreover, if \(f\) is a strict morphism of \(G\) into \(G'\) , then \(g\) is a strict morphism of \(G / H\) onto \(f(G) / f(H)\) ; for if \(U\) is open in \(G / H\) , and if \(\varphi\) (resp. \(\varphi'\) ) denotes the canonical mapping of \(G\) onto \(G / H\) {[}resp. of \(f(G)\) onto \(f(G) / f(H)\) {]}, then we have \(g(u) = \varphi'(f(\varphi^{-1}(u)))\) , and since \(\varphi^{-1}(u)\) is open in \(G\) , it follows that \(g(u)\) is open in \(f(G) / f(H)\) , which proves our assertion.
\end{enumerate}

